\section{Compact superlevels for almost every target}\label{sec:escape}

The affine law must next control escape in the source half-plane.
Our objective is compactness of positive superlevels, which will make
the subsequent critical-point counts finite above each positive value.
For \(F\in\KK\) and \(\xi\in\C\), put
\[
 V_\xi^F(x+iy)=\frac{y^k}{|F(x+iy)-\xi|},
\]
with value \(+\infty\) at the possible pole.
The choice \(3k=\beta+2\) lets a submean estimate for \(|1/F|^3\)
convert the weighted area tail below into decay of \(V_0^F\) at
source infinity.
We write \(dA\) for planar Lebesgue area.

\begin{proposition}[Compact positive superlevels]\label{prop:compact-levels}
For \(P\)-almost every \(F\), and then for area-almost every
\(\xi\in\C\), every closed positive superlevel
\(\{V_\xi^F\ge\delta\}\), \(\delta>0\), is compact in \(\HH\).
\end{proposition}

We prove this by controlling infinity and then the real boundary.
The order of these two controls is useful: the first supplies
a horizontal preimage-length bound for the second. That bound is
uniform in the height, with the map and compact target set fixed.

\begin{lemma}[The weighted area tail]\label{lem:area-tail}
Writing \(S_0=\{z\in\HH:|m(z)|\ge1/2\}\), almost every \(F\) satisfies
\begin{equation}\label{eq:area-tail}
 I(F):=\int_{S_0}y^\beta|F(z)|^{-3}\,dA(z)<\infty.
\end{equation}
\end{lemma}
\begin{proof}
For fixed \(z\ne i\), apply \eqref{eq:law} to
\(\Phi(F)=|F(z^\#)|^{-3}\).
The identity
\((T_zF)(z^\#)=-F(z)Q_F(z)/y\) gives
\[
 \E\bigl[|F'(z)|^2|F(z)|^{-3}\bigr]
 =y^{-\beta-3}\E|F(z^\#)|^{-3}.
\]
The integral of the left side over \(S_0\) is uniformly finite:
by injectivity and \eqref{eq:strip}, it is bounded by
\(\int_{|\omega|\ge c}|\omega|^{-3}\,dA(\omega)\).
The involution \(z\mapsto z^\#\) preserves \(|m(z)|\), inverts
the height, and has absolute Jacobian \(y^{-3}\).
Changing variables on the right leaves exactly the weight \(y^\beta\).
Tonelli proves finite expectation of \(I\), hence the assertion.
\end{proof}

\begin{lemma}[Escape at infinity and horizontal length]\label{lem:infinity}
If \(I(F)<\infty\), then
\begin{equation}\label{eq:infinity}
 \frac{y^k}{|F(x+iy)|}\longrightarrow0
 \qquad\text{as } |x|+y\longrightarrow\infty.
\end{equation}
For every compact target set \(K\), its preimage restricted to
\(y\ge\epsilon>0\) is compact, its whole preimage has bounded height,
and
\begin{equation}\label{eq:horizontal}
 \sup_{s>0}\int_\R\ind_K(F(x+is))\,dx<\infty.
\end{equation}
\end{lemma}
\begin{proof}
If \(w=x'+iy'\in B(z,y/2)\), then \(y/2\le y'\le3y/2\) and
\(|x|+y\le |x'|+3y'\).
Thus these disks escape every bounded \(|x'|+y'\) region as
their centers do. The complement of \(S_0\) is bounded, so
the disks are eventually contained in \(S_0\) and avoid the zero of \(F\).
Submean for \(|1/F|^3\), together with \(3k=\beta+2\), gives
\[
 \left(\frac{y^k}{|F(z)|}\right)^3
 \le C_\beta\int_{B(z,y/2)}
      (\im w)^\beta|F(w)|^{-3}\,dA(w).
\]
The right side tends to zero as a tail of \eqref{eq:area-tail}.
For \(|F(z)|\) bounded and \(y\ge\epsilon\), this excludes source
infinity, proving strip compactness. It also excludes unbounded
heights in the whole preimage of \(K\).

Choose a nonnegative smooth compactly supported \(\psi\ge\ind_K\)
and put \(H(y)=\int_\R\psi(F(x+iy))\,dx\).
For fixed \(s>0\), the support in \(y\ge s/2\) is compact by the
strip assertion. Hence differentiation and integration by parts give
\begin{align*}
 H(s)&=\int_s^\infty(y-s)H''(y)\,dy\\
 &=\int_{y>s}(y-s)|F'(z)|^2(\Delta\psi)(F(z))\,dA(z).
\end{align*}
Here the integral of \(\partial_{xx}(\psi\circ F)\) is zero and
\(\Delta(\psi\circ F)=|F'|^2(\Delta\psi)\circ F\).
If \(Y_\psi\) bounds the preimage heights of \(\supp\psi\), the
absolute value is at most
\[
 Y_\psi\int_\HH |F'(z)|^2|(\Delta\psi)(F(z))|\,dA(z)
 \le Y_\psi\int_\C|\Delta\psi|\,dA.
\]
The last inequality uses injectivity, and the bound is independent
of \(s\). This proves \eqref{eq:horizontal}.
\end{proof}

\begin{proof}[Proof of Proposition~\ref{prop:compact-levels}]
Fix \(F\) satisfying Lemma~\ref{lem:area-tail}. The preceding
lemma controls source infinity. It remains to control approach to the
real boundary, uniformly in the horizontal coordinate for almost every
target.
For omitted targets, meaning \(\xi\notin F(\HH)\), restrict to
\(|\xi|\le R\) and fix \(\delta>0\).
At scale \(h=2^{-\ell}\), \(\ell\ge0\), partition the line into
cells \([jh,(j+1)h)\). Mark a cell if some point \(z=x+iy\)
with \(x\) in that cell and \(h\le y\le2h\) satisfies
\(|F(z)-\xi|\le y^k/\delta\) for at least one such omitted target.
At a witness, \eqref{eq:local-quarter} on \(B(z,y)\) gives
\(y|F'(z)|\le4|F(z)-\xi|\).
For every \(x'\) in the same cell, the rebased point
\((x'-x)/y+i h/y\) belongs to the compact rectangle
\([-1,1]+i[1/2,1]\). Lemma~\ref{lem:compact} applied to \(T_zF\)
therefore gives
\begin{equation}\label{eq:cell-cover}
 |F(x'+ih)-\xi|\le C_\delta h^k.
\end{equation}
This bound holds for every witness and its target.

Each marked cell's entire segment at height \(h\) is mapped into
\(\{|\omega|\le R+C_\delta\}\).
Equation~\eqref{eq:horizontal} thus bounds the number of marked
cells by \(C/h\).
Disks of radius \(C_\delta h^k\) centered at their midpoint images
cover all targets witnessed at this scale, and have summed area
\(O(h^{2k-1})\).
The sum over dyadic scales is finite because \(k>1/2\).
The limsup of these measurable covering unions has area zero.
Taking integer \(R>0\) and \(\delta=1/n\), \(n\ge1\), shows that
for almost every omitted target,
\[
 \sup_x V_\xi^F(x+iy)\longrightarrow0\qquad(y\downarrow0).
\]
Indeed a failure at arbitrarily small heights gives witnesses
in infinitely many dyadic bands for one such threshold.

If \(\xi\in F(\HH)\), a small closed disk about \(\xi\) contained
in \(F(\HH)\) has compact inverse image under the homeomorphism
onto the image. Below its minimum preimage height,
\(|F(z)-\xi|\) is uniformly positive, giving the same limit.
At heights bounded below, \eqref{eq:infinity} gives
\(|F(z)|\to\infty\) and \(V_\xi^F(z)\to0\) at source infinity.
Finally \(\{V_\xi^F\ge\delta\}\) is relatively closed, since it is
\(\{|F(z)-\xi|\le y^k/\delta\}\), also at a pole.
The two escape exclusions prove its compactness.
\end{proof}
